\documentclass{zkdl-presentation-template}
\usetikzlibrary{decorations.pathmorphing}

\title[Code-based Cryptography]{\textbf{Introduction to Code-based Cryptography}}
\author{Blockstream Research}
\date{\today}
\homepage{research.blockstream.com}
\github{BlockstreamResearch/zk-book}

\begin{document}
    \maketitleframe

    \begin{frame}{Plan}
      \tableofcontents[hideallsubsections]
    \end{frame}

    \begin{frame}{What will we learn today?}
        \begin{itemize}
            \item Error-correcting codes: encode $k$ symbols into $n > k$ symbols so that any two codewords differ in many positions.\pause
            \item Codes are everywhere in cryptography: Shamir's secret sharing \textbf{is} a Reed--Solomon code, McEliece encryption, post-quantum candidates.\pause
            \item In proof systems: a Merkle tree gives only \textbf{local access} to a committed vector. Committing to an \textbf{encoding} with large distance makes local access sufficient.\pause
            \item This is the principle behind STARKs, Ligero, Brakedown, Binius: only hash functions, no trusted setup, plausibly post-quantum.
        \end{itemize}
        \pause

        \begin{remarkblock}{Goal}
            Linear codes, Reed--Solomon codes $\mathsf{RS}[\mathbb{F},\Omega,k]$, and the notion of \textbf{proximity} to a code --- the central tool behind STARKs.
        \end{remarkblock}
    \end{frame}

    \section{Error-Correcting Codes}

    \subsection{Motivation}

    \begin{frame}{Motivation}
        Fix an \textbf{alphabet} $\Sigma$ (say, $\Sigma = \{0,1\}$ for binary strings).

        \begin{alertblock}{Problem}
            Alice wants to send Bob a message $m \in \Sigma^k$. Alice sits in Germany, while Bob is on the Moon. Whatever message Alice picks, two random letters of $m$ get corrupted. How can Alice transmit the exact message $m$ to Bob?
        \end{alertblock}
        \pause

        \textbf{Idea:}
        \begin{itemize}
            \item An \textbf{encoder} maps $k$ symbols to $n > k$ symbols, called a \textbf{codeword}: we add $n-k$ redundant symbols.\pause
            \item The receiver gets a \textbf{received word} of $n$ symbols, which may differ from the codeword in some positions.\pause
            \item A \textbf{decoder} recovers $m$ using the redundant information.
        \end{itemize}
    \end{frame}

    \begin{frame}{Transmission over a Noisy Channel}
        \begin{figure}
            \centering
            \resizebox{\linewidth}{!}{%
            \begin{tikzpicture}[
                cell/.style={draw=oc-gray-7, thick, minimum size=0.5cm, inner sep=0pt, font=\small\ttfamily},
                msg/.style={cell, fill=oc-blue-1},
                rdn/.style={cell, fill=oc-teal-1},
                err/.style={cell, fill=oc-red-2, draw=oc-red-8, very thick},
                proc/.style={draw=oc-gray-7, thick, rounded corners, fill=oc-gray-1, minimum height=0.8cm, minimum width=1.7cm, font=\small},
                arr/.style={-stealth, very thick, oc-gray-7},
                lab/.style={font=\footnotesize},
                brace/.style={decorate, decoration={brace, amplitude=4pt, mirror}, thick, oc-gray-7},
            ]
                % top row
                \foreach \b [count=\i from 0] in {1,0,1,1} {
                    \node[msg] (m\i) at (0.5*\i, 0) {\b};
                }
                \node[lab, above=2pt] at (0.75, 0.25) {message $m \in \Sigma^k$};
                \node[proc] (enc) at (3.9, 0) {Encoder};
                \foreach \b [count=\i from 0] in {1,0,1,1} {
                    \node[msg] (c\i) at (6.4+0.5*\i, 0) {\b};
                }
                \foreach \b [count=\i from 4] in {0,1,0} {
                    \node[rdn] (c\i) at (6.4+0.5*\i, 0) {\b};
                }
                \node[lab, above=2pt] at (7.9, 0.25) {codeword $c \in \Sigma^n$};
                \draw[brace] (8.15, -0.32) -- node[lab, below=4pt] {$n-k$ redundant symbols} (9.65, -0.32);
                \draw[arr] (m3.east) -- (enc.west);
                \draw[arr] (enc.east) -- (c0.west);

                % noisy channel
                \uncover<2->{
                \draw[-stealth, very thick, oc-red-7, decorate, decoration={snake, amplitude=1.5pt, segment length=7pt, post length=4pt}]
                    (7.4, -1.05) -- node[right=4pt, lab, text=oc-red-8] {noisy channel} (7.4, -2.05);
                \foreach \b [count=\i from 0] in {1,0,0,1} {
                    \ifnum\i=2
                        \node[err] (y\i) at (6.4+0.5*\i, -2.6) {\b};
                    \else
                        \node[msg] (y\i) at (6.4+0.5*\i, -2.6) {\b};
                    \fi
                }
                \foreach \b [count=\i from 4] in {0,1,0} {
                    \node[rdn] (y\i) at (6.4+0.5*\i, -2.6) {\b};
                }
                \node[lab, below=2pt] at (7.9, -2.85) {received word $y \in \Sigma^n$};
                \draw[-stealth, thick, oc-red-8] (6.45, -1.75) node[above, lab, text=oc-red-8, inner sep=1pt] {error} -- (y2.north west);
                }

                % decoder
                \uncover<3->{
                \node[proc] (dec) at (3.9, -2.6) {Decoder};
                \foreach \b [count=\i from 0] in {1,0,1,1} {
                    \node[msg] (r\i) at (0.5*\i, -2.6) {\b};
                }
                \node[lab, below=2pt] at (0.75, -2.85) {recovered message $m$};
                \draw[arr] (y0.west) -- (dec.east);
                \draw[arr] (dec.west) -- (r3.east);
                \draw[densely dashed, thick, oc-blue-7] (0.75, -0.3) -- node[right, lab, text=oc-blue-8] {same} (0.75, -2.3);
                }
            \end{tikzpicture}%
            }
            \caption{The encoder expands $k = 4$ symbols into $n = 7$ symbols. The channel corrupts one symbol, and the decoder still recovers the original message.}
        \end{figure}
    \end{frame}

    \begin{frame}{Code}
        \begin{definition}[Code]
            A \textbf{code} of block length $n$ over an alphabet $\Sigma$ is a subset $\code \subseteq \Sigma^n$, whose elements are called \textbf{codewords}.
            \begin{itemize}
                \item $\#\code$ is the number of codewords;
                \item $k := \log_{\#\Sigma} \#\code$ is the number of message symbols.
            \end{itemize}
            An \textbf{encoding} for $\code$ is an injective map $\Sigma^k \to \Sigma^n$ whose image is $\code$.
        \end{definition}
        \pause

        \begin{questionblock}{Question}
            How do we measure ``how much'' two words differ?
        \end{questionblock}
    \end{frame}

    \subsection{Hamming Metric}

    \begin{frame}{Hamming Distance}
        \begin{definition}[Hamming distance]
            Given $\mathbf{x}, \mathbf{y} \in \Sigma^n$, the \textbf{Hamming distance} $\Delta(\mathbf{x}, \mathbf{y})$ is the number of positions in which they differ:
            \begin{equation*}
                \Delta(\mathbf{x}, \mathbf{y}) := \#\{i \in [n]: x_i \neq y_i\}.
            \end{equation*}
        \end{definition}
        \pause

        \begin{example}
            Let $\Sigma := \{0,1\}$, $\mathbf{x} := 00000$ and $\mathbf{y} := 10100$. Then $\Delta(\mathbf{x}, \mathbf{y}) = 2$: they differ on positions $0$ and $2$.
        \end{example}
        \pause

        \begin{questionblock}{Question \#1}
            Check that $(\Sigma^n,\Delta)$ is a metric space: $\Delta(\mathbf{x},\mathbf{x}) = 0$, positivity, symmetry, and the triangle inequality $\Delta(\mathbf{x},\mathbf{z}) \leq \Delta(\mathbf{x},\mathbf{y}) + \Delta(\mathbf{y},\mathbf{z})$.
        \end{questionblock}
    \end{frame}

    \begin{frame}{Relative Hamming Distance}
        Hamming distance alone does not give the full picture: two vectors of length $1000$ that differ in $5$ places are much more similar than two vectors of length $10$ that differ in $5$ places.
        \pause

        \begin{definition}[Relative Hamming distance]
            For $\mathbf{x}, \mathbf{y} \in \Sigma^n$, the \textbf{relative Hamming distance} is
            \begin{equation*}
                \delta(\mathbf{x}, \mathbf{y}) := \frac{1}{n} \Delta(\mathbf{x}, \mathbf{y}) \in [0,1].
            \end{equation*}
        \end{definition}
        \pause

        \begin{questionblock}{Question \#2}
            Over $\mathbb{F}_3$, let $\mathbf{x} = (1,0,2,2,1)$ and $\mathbf{y} = (1,2,2,0,1)$. What are $\Delta(\mathbf{x},\mathbf{y})$ and $\delta(\mathbf{x},\mathbf{y})$?
        \end{questionblock}
    \end{frame}

    \subsection{Minimum Distance}

    \begin{frame}{Minimum Distance}
        A code's effectiveness depends on the whole set: the worst-case spacing between two codewords is the key metric.

        \begin{definition}[Minimum distance]
            The \textbf{minimum distance} of a code $\code \subseteq \Sigma^n$ is
            \begin{equation*}
                d_{\min}(\code) := \min_{\mathbf{x}\neq \mathbf{y} \in \code} \Delta(\mathbf{x}, \mathbf{y}),
            \end{equation*}
            and the \textbf{relative minimum distance} is $\mu(\code) := \frac{1}{n} d_{\min}(\code)$.
        \end{definition}
        \pause

        If $d_{\min}(\code) = d$, then:
        \begin{itemize}
            \item fewer than $d$ errors never produce another valid codeword;\pause
            \item with fewer than $d/2$ errors, the corrupted word is strictly closer to the original codeword than to any other.
        \end{itemize}
    \end{frame}

    \begin{frame}{Unique Decoding}
        \begin{proposition}
            Let $\code \subseteq \Sigma^n$ be a code. For every $\mathbf{w} \in \Sigma^n$ there exists at most one codeword $\mathbf{x} \in \code$ such that $\delta(\mathbf{x},\mathbf{w}) < \mu(\code)/2$.
        \end{proposition}
        \pause

        \textbf{Proof.} If $\mathbf{x} \neq \mathbf{x}'$ were both that close, then
        \begin{equation*}
            \delta(\mathbf{x}, \mathbf{x}') \leq \delta(\mathbf{x}, \mathbf{w}) + \delta(\mathbf{w}, \mathbf{x}') < \mu(\code),
        \end{equation*}
        which contradicts the definition of $\mu(\code)$. \qed
        \pause

        \begin{figure}
            \centering
            \begin{tikzpicture}[
                scale=0.68, every node/.style={scale=0.85},
                ball/.style={draw=oc-blue-6, thick, densely dashed, fill=oc-blue-1},
                cw/.style={circle, fill=oc-blue-8, inner sep=1.6pt},
                word/.style={circle, fill=oc-red-7, inner sep=1.6pt},
                lab/.style={font=\footnotesize},
                tag/.style={lab, fill=white, rounded corners=1pt, inner sep=1.5pt},
            ]
                \def\r{2}
                \coordinate (x) at (0, 0);
                \coordinate (y) at (4, 0);
                \coordinate (w) at (0.5, 1.2);
                \draw[ball] (x) circle (\r);
                \draw[ball] (y) circle (\r);
                \draw[very thick, oc-gray-8] (x) -- node[tag, below=3pt] {$\delta(\mathbf{x}, \mathbf{x}') \geq \mu(\code)$} (y);
                \draw[very thick, oc-teal-7] (x) -- node[tag, left=2pt, text=oc-teal-9, fill=none] {$< \mu(\code)/2$} (w);
                \draw[very thick, oc-red-7, densely dotted] (w) -- node[tag, above=3pt, sloped, text=oc-red-8, pos=0.74] {$> \mu(\code)/2$} (y);
                \draw[thick, oc-blue-7, -stealth] (y) -- node[lab, below, sloped, text=oc-blue-8] {$\mu(\code)/2$} ++(-60:\r);
                \node[cw, label={[lab]below left:$\mathbf{x}$}] at (x) {};
                \node[cw, label={[lab]above right:$\mathbf{x}'$}] at (y) {};
                \node[word, label={[lab]above:$\mathbf{w}$}] at (w) {};
            \end{tikzpicture}
        \end{figure}
    \end{frame}

    \begin{frame}{Unique Decoding Radius}
        \begin{figure}
            \centering
            \resizebox{0.95\linewidth}{!}{%
            \begin{tikzpicture}[
                ball/.style={draw=oc-gray-5, thick, densely dashed, fill=oc-gray-1},
                cw/.style={circle, fill=oc-gray-7, inner sep=1.3pt},
                src/.style={circle, fill=oc-blue-8, inner sep=1.6pt},
                word/.style={circle, fill=oc-red-7, inner sep=1.6pt},
                lab/.style={font=\footnotesize},
            ]
                \newcommand{\codewords}{
                    \foreach \name/\px/\py in {n1/1.3/0.35, n2/-0.9/1.5, n3/-1.7/-0.4, n4/0.3/-1.75, n5/1.9/-1.4, n6/1.6/1.75} {
                        \coordinate (\name) at (\px, \py);
                        \draw[ball] (\name) circle (0.6);
                        \node[cw] at (\name) {};
                    }
                    \coordinate (x) at (0,0);
                    \draw[draw=oc-blue-6, thick, densely dashed, fill=oc-blue-1] (x) circle (0.6);
                }
                % (a) few errors
                \begin{scope}[scale=1.1]
                    \codewords
                    \fill[oc-cyan-2, opacity=0.25] (x) circle (0.5);
                    \draw[oc-cyan-6, thick] (x) circle (0.5);
                    \coordinate (w) at (-0.33, 0.33);
                    \draw[very thick, oc-red-7, -stealth, shorten >=1.5pt] (x) -- (w);
                    \node[src, label={[lab, label distance=-2pt]below right:$\mathbf{x}$}] at (x) {};
                    \node[word, label={[lab, label distance=-2pt]above left:$\mathbf{w}$}] at (w) {};
                    \draw[thick, oc-gray-8, -stealth] (n5) -- ++(0:0.6);
                    \node[lab, below] at (1.9, -2.0) {$\mu(\code)/2$};
                    \node[lab] at (0, -2.95) {(a) fewer than $\mu(\code)/2$ errors};
                \end{scope}
                % (b) many errors
                \uncover<2->{
                \begin{scope}[xshift=5.9cm, scale=1.1]
                    \codewords
                    \foreach \n in {n1, n2, n3, n4} {
                        \draw[draw=oc-red-5, thick, densely dashed, fill=oc-red-1] (\n) circle (0.6);
                        \node[cw] at (\n) {};
                    }
                    \fill[oc-red-4, opacity=0.25] (x) circle (1.25);
                    \draw[oc-red-7, thick] (x) circle (1.25);
                    \coordinate (w) at (0.96, 0.26);
                    \draw[very thick, oc-red-7, -stealth, shorten >=1.5pt] (x) -- (w);
                    \node[src, label={[lab]below left:$\mathbf{x}$}] at (x) {};
                    \node[word, label={[lab, label distance=-2pt]below:$\mathbf{w}$}] at (w) {};
                    \node[lab, right=1pt] at (n1) {$\mathbf{x}'$};
                    \node[lab] at (0, -2.95) {(b) more than $\mu(\code)/2$ errors};
                \end{scope}
                }
            \end{tikzpicture}%
            }
        \end{figure}
        \vspace{-5pt}
        \begin{itemize}
            \item<1-> Balls of relative radius $\mu(\code)/2$ around codewords never overlap. In (a), $\mathbf{w}$ stays in the ball of $\mathbf{x}$ and decodes to $\mathbf{x}$.
            \item<2-> In (b), $\mathbf{w}$ is closer to $\mathbf{x}'$: the decoder returns the wrong codeword.
        \end{itemize}
        \uncover<3->{$\mu(\code)/2$ is called the \textbf{unique decoding radius} of $\code$.}
    \end{frame}

    \begin{frame}{Example: Repetition Code}
        \begin{example}[Repetition code]
            The repetition code over $\mathbb{F}_2$ of block length $3$ repeats a bit three times: $\code=\{000, 111\}$. Here $k=1$ and $d_{\min}(\code)=3$.
        \end{example}
        \pause

        \begin{itemize}
            \item \cmark\ One bit flipped, $000 \mapsto 100$: still closer to $000$ than to $111$, since $d/2 = 3/2 > 1$.\pause
            \item \xmark\ Two bits flipped, $000 \mapsto 110$: closer to $111$, decoding fails, since two errors are not fewer than $d/2$.
        \end{itemize}
    \end{frame}

    \subsection{Linear Codes}

    \begin{frame}{Hamming Weight}
        Computing $d_{\min}(\code)$ from the definition compares every pair. Practical codes offer a shortcut.
        \pause

        \begin{definition}[Hamming weight]
            The \textbf{Hamming weight} $\|\mathbf{x}\|_0$ of $\mathbf{x} \in \Sigma^n$ is the number of non-zero symbols in $\mathbf{x}$.
        \end{definition}
        \pause

        When $\Sigma$ is a field,
        \begin{equation*}
            \Delta(\mathbf{u}, \mathbf{v}) = \|\mathbf{u} - \mathbf{v}\|_0,
        \end{equation*}
        that is, two vectors differ exactly in the coordinates where their difference is non-zero.
    \end{frame}

    \begin{frame}{Linear Codes}
        Cryptographic codes are never arbitrary subsets of $\Sigma^n$: they have algebraic structure.

        \begin{definition}[Linear code]
            Let $\mathbb{F}_q$ be a finite field. A code $\code \subseteq \mathbb{F}_q^n$ is a \textbf{linear code} if it is a linear subspace of $\mathbb{F}_q^n$. If $\code$ has dimension $k$ and minimum distance $d$, we call it an $[n,k,d]_q$ code.
        \end{definition}
        \pause

        \begin{questionblock}{Question \#3}
            Show that any linear encoding $\mathbb{F}^k \to \code$ is $\mathbf{x} \mapsto \mathbf{G}\mathbf{x}$ for some $\mathbf{G} \in \mathbb{F}^{n \times k}$. What is $\text{rank}(\mathbf{G})$?
        \end{questionblock}
    \end{frame}

    \begin{frame}{Distance Equals Weight}
        The difference of two codewords of a linear code is again a codeword, so the minimum over pairs collapses into a minimum over single codewords.

        \begin{proposition}
            Let $\code \subseteq \mathbb{F}^n$ be a linear code. Then
            \begin{equation*}
                \mu(\code) = \frac{1}{n}\cdot \min_{0 \neq \mathbf{x} \in \code} \|\mathbf{x}\|_0.
            \end{equation*}
        \end{proposition}
        \pause

        \textbf{Proof idea.}
        \begin{itemize}
            \item[$\leq$] $0 \in \code$, and $\Delta(0, \mathbf{x}) = \|\mathbf{x}\|_0$.\pause
            \item[$\geq$] If $\Delta(\mathbf{x}, \mathbf{y}) = d_{\min}(\code)$, then $0 \neq \mathbf{x} - \mathbf{y} \in \code$ and $\|\mathbf{x} - \mathbf{y}\|_0 = \Delta(\mathbf{x}, \mathbf{y})$. \qed
        \end{itemize}
    \end{frame}

    \begin{frame}{Rate}
        \begin{definition}[Rate]
            Let $\code \subseteq \mathbb{F}^n$ be a linear code of dimension $k$. The \textbf{rate} of $\code$ is $\rho(\code) := k/n$.
        \end{definition}
        \pause

        \begin{itemize}
            \item A codeword has $n$ symbols but carries only $k$ symbols of information: $\rho$ is the fraction of actual payload.\pause
            \item Rate and distance pull in opposite directions: every symbol is spent either on carrying the message or on keeping codewords apart.
        \end{itemize}
    \end{frame}

    \begin{frame}{Example: Parity Code}
        \begin{example}[Parity code]
            The \textbf{parity code} consists of all vectors with an even number of ones:
            \begin{equation*}
                \code = \left\{ \mathbf{c} \in \mathbb{F}_2^n : c_1 + c_2 + \dots + c_n = 0 \right\}.
            \end{equation*}
            To encode $\mathbf{m} \in \mathbb{F}_2^{n-1}$, append a \textbf{parity bit}. For $n = 4$: $101 \mapsto 1010$, $111 \mapsto 1111$.
        \end{example}
        \pause

        \begin{itemize}
            \item Linear, of dimension $n-1$; weight $1$ is impossible and $110\ldots0 \in \code$, so $d_{\min}(\code) = 2$.\pause
            \item Hence $\code$ is an $[n, n-1, 2]_2$ code with rate $\rho(\code) = 1 - 1/n$.\pause
            \item A single error is \textbf{detected}, but not \textbf{corrected}: $1110$ is at distance $1$ from both $1010$ and $1111$.
        \end{itemize}
    \end{frame}

    \begin{frame}{Singleton Bound}
        To guarantee distance $d$, a code must dedicate at least $d-1$ of its $n$ symbols to redundancy.

        \begin{proposition}[Singleton bound]
            Every $[n,k,d]_q$ code satisfies $k \leq n-d+1$, equivalently $\#\code \leq q^{n-d+1}$.
        \end{proposition}
        \pause

        \textbf{Proof.} Delete the first $d-1$ symbols of every codeword.
        \begin{itemize}
            \item Distinct codewords differ in $\geq d$ positions, so they still differ in one of the remaining $n-d+1$ positions.\pause
            \item There are at most $q^{n-d+1}$ such strings, so $q^k = \#\code \leq q^{n-d+1}$. \qed
        \end{itemize}
    \end{frame}

    \begin{frame}{MDS Codes}
        \begin{definition}[MDS code]
            An $[n,k,d]_q$ code is \textbf{maximum distance separable} (MDS) if $d = n-k+1$, that is, if it attains the Singleton bound.
        \end{definition}
        \begin{itemize}
            \item No redundant symbol is wasted.
            \item Any $k$ symbols of a codeword determine the whole codeword: the message survives any $n-k$ erasures.
        \end{itemize}
        \pause

        \begin{example}
            The repetition code $[n,1,n]_2$ and the parity code $[n,n-1,2]_2$ are MDS. A $[7,4,3]_2$ code is not MDS, since $n-k+1 = 4 > 3$.
        \end{example}
        \pause

        \begin{questionblock}{Question \#4}
            Which of $[7,4,4]_q$, $[7,4,5]_q$, $[5,1,5]_q$, $[6,3,4]_q$ cannot be achieved by any linear code?
        \end{questionblock}
    \end{frame}

    \section{Reed--Solomon Codes}

    \begin{frame}{Idea: Encode into Polynomials}
        As with many ZK protocols, encode data into polynomials and store their \textbf{evaluations over a large domain}.
        \pause
        \begin{algoen}
            \item Put the message $(m_i)_{i \in [k]} \in \mathbb{F}^k$ into the coefficients: $f := \sum_{i \in [k]}m_iT^i \in \mathbb{F}[T]^{<k}$.\pause
            \item Fix an \textbf{evaluation domain} $\Omega = (\omega_i)_{i \in [n]}$ with $n > k$.\pause
            \item The codeword is $(f(\omega_i))_{i \in [n]}$.
        \end{algoen}
        \pause

        \begin{remarkblock}{Recall: Shamir's secret sharing}
            Decoding follows the same idea: if few evaluations are corrupted, we can still interpolate the polynomial $f$ of degree $k-1$.
        \end{remarkblock}
    \end{frame}

    \begin{frame}{Reed--Solomon Code}
        \textbf{Restriction:} for $f \in \mathbb{F}[T]^{<k}$ and a domain $\Omega \subseteq \mathbb{F}$, we write $\overline{f}: \Omega \to \mathbb{F}$ for $f$ viewed only on inputs from $\Omega$.

        \begin{remark}
            Feel free to treat $\overline{f}$ as the vector $(f(\omega_0),\dots,f(\omega_{n-1}))$ for $\Omega := (\omega_0,\dots,\omega_{n-1})$.
        \end{remark}
        \pause

        \begin{definition}[Reed--Solomon code]
            The \textbf{Reed--Solomon code} over $\mathbb{F}$, evaluation domain $\Omega \subseteq \mathbb{F}$ (with $\#\Omega = n > k$), and degree $k$, is the linear code
            \begin{equation*}
                \mathsf{RS}[\mathbb{F},\Omega,k] := \{\overline{f}: \Omega \to \mathbb{F}, \; f \in \mathbb{F}[T]^{<k}\}.
            \end{equation*}
        \end{definition}
    \end{frame}

    \begin{frame}{Example: $\mathsf{RS}[\mathbb{F}_7, \mathbb{F}_7, 3]$}
        \begin{columns}[T]
            \begin{column}{0.47\textwidth}
                \small
                Take $\Omega = \mathbb{F}_7$, $n = 7$, $k = 3$.
                \begin{itemize}
                    \item $(1, 2, 3) \mapsto f = 1 + 2T + 3T^2$,
                    \item $(4, 0, 2) \mapsto g = 4 + 2T^2$.
                \end{itemize}
                \pause
                Codewords:
                \begin{align*}
                    \overline{f} &= (1, 6, 3, 6, 1, 2, 2), \\
                    \overline{g} &= (4, 6, 5, 1, 1, 5, 6).
                \end{align*}
                \pause
                They agree only at $1$ and $4$: $f - g = (T-1)(T-4)$ has at most $k - 1 = 2$ roots, so $\overline{f}, \overline{g}$ differ in $5$ of $7$ positions.
            \end{column}
            \begin{column}{0.5\textwidth}
                \centering
                \resizebox{\linewidth}{!}{%
                \begin{tikzpicture}[
                    x=0.8cm, y=0.5cm,
                    lab/.style={font=\footnotesize},
                    cell/.style={draw=oc-gray-7, thick, minimum width=0.8cm, minimum height=0.5cm, inner sep=0pt, font=\small\ttfamily},
                ]
                    \foreach \v in {0,...,6} {
                        \draw[oc-gray-3] (\v, 0) -- (\v, 6);
                        \draw[oc-gray-3] (0, \v) -- (6, \v);
                        \node[lab, left=3pt] at (0, \v) {\v};
                    }
                    \node[lab, rotate=90] at (-1, 3) {value in $\mathbb{F}_7$};
                    \uncover<3->{
                    \foreach \v in {1, 4} {
                        \fill[oc-orange-2, opacity=0.6] (\v-0.5, -0.5) rectangle (\v+0.5, 6.5);
                    }}
                    \uncover<2->{
                    \foreach \px/\py in {0/4, 1/6, 2/5, 3/1, 4/1, 5/5, 6/6} {
                        \node[rectangle, fill=oc-teal-6, inner sep=2.6pt] at (\px, \py) {};
                    }
                    \foreach \px/\py in {0/1, 1/6, 2/3, 3/6, 4/1, 5/2, 6/2} {
                        \node[circle, fill=oc-blue-7, inner sep=1.7pt] at (\px, \py) {};
                    }
                    \foreach \v [count=\i from 0] in {1,6,3,6,1,2,2} {
                        \node[cell, fill=oc-blue-1] at (\i, -1.6) {\v};
                    }
                    \foreach \v [count=\i from 0] in {4,6,5,1,1,5,6} {
                        \node[cell, fill=oc-teal-1] at (\i, -2.6) {\v};
                    }}
                    \uncover<3->{
                    \foreach \v in {1, 4} {
                        \draw[oc-orange-7, very thick] (\v-0.5, -3.1) rectangle (\v+0.5, -1.1);
                    }}
                    \node[lab, left=3pt, text=oc-blue-8] at (-0.5, -1.6) {$\overline{f}$};
                    \node[lab, left=3pt, text=oc-teal-8] at (-0.5, -2.6) {$\overline{g}$};
                    \foreach \v in {0,...,6} {
                        \node[lab] at (\v, -3.7) {\v};
                    }
                    \node[lab] at (3, -4.7) {evaluation point $\omega \in \Omega$};
                \end{tikzpicture}%
                }
            \end{column}
        \end{columns}
    \end{frame}

    \begin{frame}{Reed--Solomon Code is Linear}
        \begin{proposition}
            $\mathsf{RS}[\mathbb{F}, \Omega, k]$ is a linear code of length $n = \#\Omega$ and dimension $k$.
        \end{proposition}
        \pause

        \textbf{Proof.}
        \begin{itemize}
            \item For $a \in \mathbb{F}$ and $f, g \in \mathbb{F}[T]^{<k}$, restriction commutes with both operations:
            \begin{equation*}
                \overline{f} + \overline{g} = \overline{f + g}, \qquad a \overline{f} = \overline{af}.
            \end{equation*}
            \pause
            \vspace{-10pt}
            \item The map $f \mapsto \overline{f}$ is linear and injective: if $\overline{f} = 0$, then $f$ has $n > k - 1$ roots, so $f = 0$. Hence the dimension is $\dim \mathbb{F}[T]^{<k} = k$. \qed
        \end{itemize}
    \end{frame}

    \begin{frame}{Reed--Solomon Codes are MDS}
        \begin{proposition}
            $\mathsf{RS}[\mathbb{F}, \Omega, k]$ with $n = \#\Omega$ and $q = \#\mathbb{F}$ is an $[n,k,n-k+1]_q$ code. That is, it attains the Singleton bound.
        \end{proposition}
        \pause

        \textbf{Proof.} Fix distinct $f, g \in \mathbb{F}[T]^{<k}$. Then
        \begin{equation*}
            \Delta(\overline{f}, \overline{g}) = n - \#\{\omega \in \Omega : f(\omega) = g(\omega)\} \geq n - (k-1),
        \end{equation*}
        since $f - g \neq 0$ has at most $k-1$ roots. The Singleton bound gives the other direction. \qed
        \pause

        \begin{remark}
            The rate is $\rho = k/n$ and $\mu = 1 - \rho + 1/n$: lowering the rate directly increases the relative distance. Since $n$ is large in practice, we use $\mu \approx 1-\rho$.
        \end{remark}
    \end{frame}

    \begin{frame}{Questions}
        \begin{questionblock}{Question \#5}
            Consider $\mathsf{RS}[\mathbb{F}_7,\Omega:=\mathbb{F}_7,2]$. What is the encoding of the message $(3,2)$?
        \end{questionblock}
        \pause

        \begin{questionblock}{Question \#6}
            Let $\code \subseteq \mathbb{F}_2^4$ be the linear code spanned by $(1,1,0,1)$ and $(0,1,1,1)$. What is $d_{\min}(\code)$?
        \end{questionblock}
    \end{frame}

    \section{Proximity to a Code}

    \subsection{List Decoding}

    \begin{frame}{Distance to a Code}
        Classical coding theory: recover the original message. Our requirement is weaker:
        \begin{itemize}
            \item the verifier does not need to know \textbf{which} polynomial the prover committed to,
            \item only whether the committed vector is close to \textbf{some} codeword.
        \end{itemize}
        \pause

        \begin{definition}[Distance to a code]
            Let $\code \subseteq \Sigma^n$ be a code and $\mathbf{w} \in \Sigma^n$. Then
            \begin{equation*}
                \Delta(\mathbf{w}, \code) := \min_{\mathbf{x} \in \code} \Delta(\mathbf{w}, \mathbf{x}), \qquad \delta(\mathbf{w}, \code) := \frac{1}{n}\Delta(\mathbf{w}, \code).
            \end{equation*}
        \end{definition}
        \pause

        $\mathbf{w}$ is \textbf{$\delta$-close} to $\code$ if $\delta(\mathbf{w}, \code) \leq \delta$, and \textbf{$\delta$-far} otherwise.
    \end{frame}

    \begin{frame}{The List of Close Codewords}
        \begin{definition}
            Let $\code$ be an $[n,k,d]$ linear code. For $\mathbf{w} \in \mathbb{F}^n$ and $\delta \in [0,1]$, define
            \begin{equation*}
                \mathsf{List}[\mathbf{w},\code,\delta] := \{\mathbf{x} \in \code: \delta(\mathbf{x}, \mathbf{w}) \leq \delta\}.
            \end{equation*}
        \end{definition}
        \pause

        In many cryptographic applications, we care about the \textbf{size} of this list.
        \pause

        \begin{example}
            By unique decoding, if $\delta < \mu(\code)/2$, then $\#\mathsf{List}[\mathbf{w},\code,\delta] \leq 1$.
        \end{example}
        \pause

        \begin{alertblock}{But\ldots}
            For $\delta \geq \mu(\code)/2$ the situation becomes very tricky, and our knowledge is quite limited.
        \end{alertblock}
    \end{frame}

    \begin{frame}{Johnson Bound}
        \begin{theorem}[Johnson bound]
            For a code $\code \subseteq \mathbb{F}^n$, fixed $\mathbf{w} \in \mathbb{F}^n$, and $0 < \delta < 1-\sqrt{1-\mu(\code)}$,
            \begin{equation*}
                \#\mathsf{List}[\mathbf{w},\code,\delta] \leq 1/\varepsilon_{\delta}, \quad \varepsilon_{\delta} := 2\sqrt{1-\mu(\code)}\,(1-\sqrt{1-\mu(\code)}-\delta).
            \end{equation*}
        \end{theorem}
        \pause

        \begin{itemize}
            \item The bound only makes sense for $\delta < 1-\sqrt{1-\mu(\code)}$: $\varepsilon_{\delta} \to 0$ as $\delta$ approaches it.\pause
            \item For $\delta \geq \mu(\code)$, in the worst case the list is large.
        \end{itemize}
        \pause

        \begin{questionblock}{Question \#7}
            Let $\code \subseteq \mathbb{F}^n$ be linear over a field of size $q$, and $\delta \geq \mu(\code)$. Show that some $\mathbf{w} \in \mathbb{F}^n$ has $\#\mathsf{List}[\mathbf{w}, \code, \delta] \geq q$.
        \end{questionblock}
    \end{frame}

    \begin{frame}{Size of the List}
        What happens between $1-\sqrt{1-\mu(\code)}$ and $\mu(\code)$? There is no general answer: it depends on the specifics of $\code$.

        \begin{figure}
            \centering
            \resizebox{\linewidth}{!}{%
            \begin{tikzpicture}[
                lab/.style={font=\footnotesize},
                seg/.style={line width=5pt},
                tick/.style={thick, oc-gray-8},
            ]
                \def\xa{0} \def\xb{2.3} \def\xc{5.6} \def\xd{8.6} \def\xe{10.8}
                \draw[seg, oc-teal-4]   (\xa, 0) -- (\xb, 0);
                \node[lab, above=5pt, text=oc-teal-9]   at ({(\xa+\xb)/2}, 0) {size $\leq 1$};
                \uncover<2->{
                \draw[seg, oc-blue-4]   (\xb, 0) -- (\xc, 0);
                \node[lab, above=5pt, text=oc-blue-9]   at ({(\xb+\xc)/2}, 0) {size $\leq 1/\varepsilon_\delta$};}
                \uncover<4->{
                \draw[seg, oc-orange-4] (\xc, 0) -- (\xd, 0);
                \node[lab, above=5pt, text=oc-orange-9] at ({(\xc+\xd)/2}, 0) {depends on $\code$};}
                \uncover<3->{
                \draw[seg, oc-red-4]    (\xd, 0) -- (\xe, 0);
                \node[lab, above=5pt, text=oc-red-9]    at ({(\xd+\xe)/2}, 0) {size $\geq q$};}
                \foreach \x/\t in {\xa/{$\delta = 0$}, \xb/{$\mu/2$}, \xc/{$1 - \sqrt{1 - \mu}$}, \xd/{$\mu$}, \xe/{$1$}} {
                    \draw[tick] (\x, -0.22) -- (\x, 0.22);
                    \node[lab, below=1pt] at (\x, -0.25) {\t};
                }
                \foreach \x/\t in {\xb/{unique decoding\\radius}, \xc/{Johnson\\bound}, \xd/{capacity\\bound}} {
                    \node[lab, below, align=center, text=oc-gray-7] at (\x, -0.75) {\t};
                }
            \end{tikzpicture}%
            }
            \caption{The size of $\mathsf{List}[\mathbf{w}, \code, \delta]$ as a function of $\delta \in [0, 1]$, where $\mu = \mu(\code)$. Adapted from Dan Boneh's ZK Whiteboard Session S2M7.}
        \end{figure}
    \end{frame}

    \subsection{List Decoding for Reed--Solomon Codes}

    \begin{frame}{List Decoding for Reed--Solomon Codes}
        For $\code = \mathsf{RS}[\mathbb{F},\Omega,k]$, substitute $\mu(\code) \approx 1-\rho(\code)$.

        \begin{proposition}
            Let $\code = \mathsf{RS}[\mathbb{F}_q,\Omega,k]$, $\mathbf{w} \in \mathbb{F}^n$, and $\rho = \rho(\code)$. Then
            \begin{enumerate}
                \item for $\delta \leq (1-\rho)/2$, $\#\mathsf{List}[\mathbf{w},\code,\delta] \leq 1$;
                \item for $\delta < 1-\sqrt{\rho}$, $\#\mathsf{List}[\mathbf{w},\code,\delta] \leq \varepsilon_{\delta}^{-1}$ for $\varepsilon_{\delta} := 2\sqrt{\rho}(1-\sqrt{\rho}-\delta)$;
                \item for $\delta \geq 1-\rho$, $\#\mathsf{List}[\mathbf{w},\code,\delta] \geq q$.
            \end{enumerate}
        \end{proposition}
        \pause

        \begin{figure}
            \centering
            \resizebox{0.95\linewidth}{!}{%
            \begin{tikzpicture}[
                lab/.style={font=\footnotesize},
                seg/.style={line width=5pt},
                tick/.style={thick, oc-gray-8},
            ]
                \def\xa{0} \def\xb{2.3} \def\xc{5.6} \def\xd{8.6} \def\xe{10.8}
                \draw[seg, oc-teal-4]   (\xa, 0) -- (\xb, 0);
                \draw[seg, oc-blue-4]   (\xb, 0) -- (\xc, 0);
                \draw[seg, oc-orange-4] (\xc, 0) -- (\xd, 0);
                \draw[seg, oc-red-4]    (\xd, 0) -- (\xe, 0);
                \node[lab, above=5pt, text=oc-teal-9]   at ({(\xa+\xb)/2}, 0) {size $\leq 1$};
                \node[lab, above=5pt, text=oc-blue-9]   at ({(\xb+\xc)/2}, 0) {size $\leq 1/\varepsilon_\delta$};
                \node[lab, above=5pt, text=oc-orange-9] at ({(\xc+\xd)/2}, 0) {unknown (only conjectured)};
                \node[lab, above=5pt, text=oc-red-9]    at ({(\xd+\xe)/2}, 0) {size $\geq q$};
                \foreach \x/\t in {\xa/{$\delta = 0$}, \xb/{$(1-\rho)/2$}, \xc/{$1 - \sqrt{\rho}$}, \xd/{$1 - \rho$}, \xe/{$1$}} {
                    \draw[tick] (\x, -0.22) -- (\x, 0.22);
                    \node[lab, below=1pt] at (\x, -0.25) {\t};
                }
                \foreach \x/\t in {\xb/{unique decoding\\radius}, \xc/{Johnson\\bound}, \xd/{capacity\\bound}} {
                    \node[lab, below, align=center, text=oc-gray-7] at (\x, -0.75) {\t};
                }
            \end{tikzpicture}%
            }
        \end{figure}
    \end{frame}

    \begin{frame}{The Proximity Prize}
        \begin{itemize}
            \item The size of the list beyond the Johnson bound is so important that the Ethereum Foundation announced the \textbf{Proximity Prize} (\href{https://proximityprize.org/}{proximityprize.org}).\pause
            \item Goal: verify or disprove the conjectures for larger $\delta$'s.\pause
            \item A comprehensive review: survey by Gal Arnon, Dan Boneh, and Giacomo Fenzi.
        \end{itemize}
        \pause

        \begin{questionblock}{Question \#8}
            Consider $\mathsf{RS}[\mathbb{F}_7,\Omega:=\mathbb{F}_7,2]$ and $\mathbf{w} := (0,1,4,2,2,4,1)$. What is $\delta(\mathbf{w}, \mathsf{RS})$?

            \textit{Hint:} $\mathbf{w}$ is the evaluation of $T^2$ over $\Omega := \mathbb{F}_7$.
        \end{questionblock}
    \end{frame}

    \subsection{Why do we need List Decoding?}

    \begin{frame}{Why do we need List Decoding?}
        Assume that $\mathbf{w}$ is $\delta$-far from $\code$.
        \begin{itemize}
            \item Whichever $\mathbf{x} \in \code$ we compare it to, they differ in more than $\delta n$ positions.\pause
            \item A uniformly random query detects a discrepancy with probability $> \delta$; $t$ independent queries miss it with probability $< (1-\delta)^t$.\pause
            \item The larger the $\delta$ we can afford, the fewer queries and the shorter the proof.
        \end{itemize}
        \pause

        \begin{alertblock}{The challenge}
            The verifier does not know in advance which codeword $\mathbf{x}$ it is checking against. An efficient test with few queries to $\mathbf{w}$ is exactly what the FRI protocol achieves.
        \end{alertblock}
    \end{frame}

    \begin{frame}{Three Regimes}
        \begin{itemize}
            \item \textbf{Unique decoding}, $\delta < (1-\rho)/2$: a query passes with probability $> (1+\rho)/2 > 1/2$, so less than one bit of security per query.\pause
            \item Beyond this radius, $\mathbf{w}$ no longer determines a single codeword: every element of $\mathsf{List}[\mathbf{w}, \code, \delta]$ is a candidate, and the soundness error grows with $\#\mathsf{List}[\mathbf{w}, \code, \delta]$.\pause
            \item \textbf{Johnson bound}, $\delta < 1 - \sqrt{\rho}$: the list is provably small, and a query passes with probability about $\sqrt{\rho}$.\pause
            \item \textbf{Capacity}, $\delta < 1 - \rho$: the list is only \textbf{conjectured} to be small; protocols using this range rely on the conjecture.
        \end{itemize}
        \pause

        \begin{example}
            For $\rho = 1/4$ and $100$ bits of security, the three regimes require roughly $148$, $100$, and $50$ queries, respectively.
        \end{example}
    \end{frame}

    \begin{frame}{Beyond Proof Systems: Keys from Biometrics}
        Two readings $\mathbf{w}, \mathbf{w}' \in \mathbb{F}_2^n$ of the same face scan are never identical, only \textbf{close} in the Hamming distance, while a key must be reproduced exactly.
        \pause

        \begin{block}{Code-offset fuzzy extractor}
            Let $\code$ be a binary $[n, k, 2t+1]_2$ code with a decoder correcting up to $t$ errors, and $\mathsf{H}$ a hash function.
            \begin{itemize}
                \item $\mathsf{Gen}(\mathbf{w})$: sample $\mathbf{c} \leftarrowS \code$; output the public \textbf{helper string} $\mathbf{h} := \mathbf{c} + \mathbf{w}$ and the key $\mathsf{sk} := \mathsf{H}(\mathbf{c})$.
                \item $\mathsf{Rep}(\mathbf{h}, \mathbf{w}')$: decode $\mathbf{h} + \mathbf{w}'$ to a codeword $\mathbf{c}'$; output $\mathsf{sk}' := \mathsf{H}(\mathbf{c}')$.
            \end{itemize}
        \end{block}
        \pause

        \begin{questionblock}{Question \#9}
            Show that $\mathsf{Rep}(\mathbf{h}, \mathbf{w}') = \mathsf{sk}$ whenever $\Delta(\mathbf{w}, \mathbf{w}') \leq t$. For a $[128, 44, 25]_2$ code, why does the key have at most $44$ bits of security, however good the biometric is?
        \end{questionblock}
    \end{frame}

    \begin{frame}{Summary}
        \begin{itemize}
            \item A code $\code \subseteq \Sigma^n$ adds redundancy; the \textbf{relative minimum distance} $\mu(\code)$ governs error correction, and $\mu(\code)/2$ is the unique decoding radius.\pause
            \item \textbf{Linear} $[n,k,d]_q$ codes: distance equals minimum weight; rate $\rho = k/n$; Singleton bound $k \leq n - d + 1$.\pause
            \item \textbf{Reed--Solomon} codes $\mathsf{RS}[\mathbb{F},\Omega,k]$: restrictions of polynomials of degree $< k$ to $\Omega$; they are MDS and $\mu \approx 1 - \rho$.\pause
            \item \textbf{Proximity}: the size of $\mathsf{List}[\mathbf{w},\code,\delta]$ is $\leq 1$ below $(1-\rho)/2$, small up to the Johnson bound $1 - \sqrt{\rho}$, and only conjectured to be small up to $1-\rho$.\pause
            \item Next: testing proximity with few queries (IOPPs and FRI).
        \end{itemize}
    \end{frame}

    \makeclosingframe
\end{document}
